\documentclass[tikz,border=0pt]{standalone}
\usepackage[T1]{fontenc}
\usepackage{helvet}
\usepackage{graphicx}
\graphicspath{{../../.local/figures/plots/}}
\renewcommand{\familydefault}{\sfdefault}
\usetikzlibrary{arrows.meta,calc}
\newcommand{\markhinge}[2]{\fill (#1,#2) circle (.27mm);}
\tikzset{
 flow/.style={-{Stealth[length=1.7mm,width=1.15mm]},line width=.75pt},
 pics/hingex/.style={code={
   \foreach \i in {-3,-2,-1,0,1,2,3}{
      \markhinge{\i}{.65*\i}\markhinge{\i}{-.65*\i}
   }
   \foreach \i in {-1,0,1}{\markhinge{\i}{.6}\markhinge{\i}{-.6}}
 }},
 pics/hingev/.style={code={
   \foreach \i in {-3,-2,-1,0,1,2,3}{
      \pgfmathsetmacro{\yy}{.5*abs(\i)+.8}
      \markhinge{\i}{\yy}\markhinge{\i}{-\yy}
   }
 }},
 pics/hingevzv/.style={code={
   \foreach \i in {-3,-2,-1,0,1,2,3}{
     \markhinge{\i}{2}\markhinge{\i}{-2}
   }
   \foreach \i in {-1.4,-.7,0,.7,1.4}{
     \markhinge{-2.5}{\i}\markhinge{2.5}{\i}
     \markhinge{-.8}{\i}\markhinge{.8}{\i}
   }
 }},
 pics/hingedxa/.style={code={
   \foreach \i in {-3,-2,-1,0,1,2,3}{
     \markhinge{\i}{.35*\i+.4}\markhinge{\i}{-.2*\i-.4}
     \markhinge{\i}{.35*\i+1.1}
   }
 }},
 pics/hingedxb/.style={code={
   \foreach \i in {-3,-2,-1,0,1,2,3}{
     \markhinge{\i}{-.3*\i+.1}\markhinge{\i}{.2*\i-.5}
     \markhinge{\i}{-.3*\i+1.0}
   }
 }}}
\tikzset{frameone/.style={line width=1.5pt,line join=miter,line cap=butt},frametwo/.style={line width=1pt,line join=miter,line cap=butt}}
\begin{document}
\begin{tikzpicture}[x=1mm,y=1mm,line width=.65pt,line cap=round,line join=miter,
 every node/.style={font=\fontsize{9}{11}\selectfont,inner sep=1pt,text=black}]
\path[use as bounding box] (0,0) rectangle (240,180);
\draw[frameone] (3,91) rectangle (78.333333,177) (82.333333,91) rectangle (157.666667,177) (161.666667,91) rectangle (237,177);
\begin{scope}[xshift=-3.333333mm]
\node at (44,172) {A. Plastic-hinge tracing};
\begin{scope}[shift={(22.15,123)}]

 \foreach \j in {0,...,7}{
   \pgfmathtruncatemacro{\phase}{mod(\j,2)}
   \pgfmathsetmacro{\xb}{6.1*\j+\phase}
   \pgfmathsetmacro{\xm}{6.1*\j+1-\phase}
   \foreach \r in {0,...,7}{
     \draw (\xb,5*\r)--(\xb,5*\r+.2)
       --(\xm,5*\r+2.3)--(\xm,5*\r+2.7)
       --(\xb,5*\r+4.8)--(\xb,5*\r+5);
   }
 }
 % Neighbouring cell columns are offset by half a vertical period.
 \foreach \c in {0,2,4,6}{
   \foreach \r in {0,...,8}{
     \draw (6.1*\c,5*\r)--({6.1*(\c+1)+1},5*\r);
   }
 }
 \foreach \c in {1,3,5}{
   \foreach \r in {0,...,7}{
     \draw (6.1*\c,5*\r+2.5)--({6.1*(\c+1)+1},5*\r+2.5);
   }
 }
 \draw[line width=.7pt] (-.6,-1.5) rectangle (44.3,0);
 \draw[line width=.7pt] (-.6,40) rectangle (44.3,40.7);
 % Mark actual wall joints; these are illustrative, not measured strains.
 \foreach \j/\r in {1/2,6/2,2/3,5/3,3/4,4/4,2/5,5/5,1/6,6/6}{
   \pgfmathtruncatemacro{\phase}{mod(\j,2)}
   \fill (6.1*\j+\phase,5*\r) circle (.5mm);
 }

\end{scope}
\draw[flow] (56,166)--(53.15,161.7);
\node at (44,118) {Finite-element deformation images};
\draw[frametwo] (32,94) rectangle (56,114);\pic[scale=2] at (44,104) {hingex};
\node at (44,110) {64 by 64 pixels};
\end{scope}
\node at (120,172) {B. K-means image clusters};
\foreach \xx/\yy/\id/\pattern in {103.5/158/C0/hingev,136.5/158/C1/hingex,103.5/140/C2/hingevzv,136.5/140/C3/hingedxa,103.5/122/C4/hingedxb}{
 \draw[frametwo] (\xx-6,\yy-6) rectangle (\xx+6,\yy+6);
 \pic[scale=1.3] at (\xx,\yy) {\pattern};\node at (\xx,\yy-9) {\id};
}
\foreach \xx/\yy/\id in {136.5/122/C5,103.5/103.5/C6,136.5/103.5/C7}{
 \draw[frametwo] (\xx-6,\yy-6) rectangle (\xx+6,\yy+6);
 \node at (\xx,\yy-9) {\id};
}
\foreach \i in {-3,...,3}{\fill (136.5+\i,122) circle (.35);}
\foreach \i in {-1,0,1}{\fill (103.5+\i,104) circle (.3);}\fill (136.5,104) circle (.3);
\begin{scope}[xshift=3.333333mm]
\node at (196,172) {C. Logistic-regression classifier};
\draw[frametwo] (168,146) rectangle (224,160);
\node[align=center] at (196,153) {KNet1\\Predict cluster label from image};
% Exact 60/20/20 split visualized proportionally; no new sample-count claim.
\node at (196,142.5) {Training / validation / testing};
\draw[frametwo] (168,134) rectangle (224,139);
\draw (201.6,134)--(201.6,139) (212.8,134)--(212.8,139);
\node at (184.8,136.5) {60\%};
\node at (207.2,136.5) {20\%};
\node at (218.4,136.5) {20\%};
\node at (196,126) {4 of 6 mechanisms identified};
\foreach \xx/\yy/\pattern/\label in {178/118/hingev/V,214/118/hingex/X,178/102/hingevzv/V+Z+V,214/102/hingedxa/Distorted X}{
 \pic[scale=1.5] at (\xx,\yy) {\pattern};\node at (\xx,\yy-6) {\label};
}
\end{scope}
\draw[flow,shorten <=.6mm,shorten >=.6mm] (78.333333,134)--(82.333333,134);\draw[flow,shorten <=.6mm,shorten >=.6mm] (157.666667,134)--(161.666667,134);
\draw[frameone] (3,23) rectangle (118,87) (122,23) rectangle (237,87);
\node[anchor=west] at (7,82) {D. Class-wise performance: eight image-cluster labels};
\node[anchor=west] at (126,82) {E. Learning-rate selection};
\node[anchor=south west,inner sep=0pt] at (10,24)
 {\input{../../.local/figures/plots/auxmech-plot.tex}};
\draw[frameone] (3,3) rectangle (237,19);
\node[anchor=west,align=left] at (7,11) {88.08 percent test accuracy against cluster-derived labels.\\Independent deformation-mode ground truth is unavailable.};
\node[anchor=east,align=right] at (233,11) {Singh, Dhari and Javanbakht (2025), 853--877.\\International Journal of Protective Structures 16(4).\\Source: Tables 1, 2 and 4; artwork: conceptual schematic.};
\end{tikzpicture}
\end{document}
